\documentclass[10pt,a4paper]{amsart}
\usepackage[margin=19mm]{geometry}
\usepackage{amsmath,amssymb,mathtools}
\usepackage[scr=rsfs,cal=euler]{mathalfa}
\usepackage{xfrac}
\usepackage[unicode,hidelinks,bookmarks=false]{hyperref}
\newcommand{\ie}{i.e.\ }
\newcommand{\cf}{cf.\ }
\newcommand{\resp}{respectively\ }
\newcommand{\Subsection}{\subsection}
\providecommand{\nobreakdash}{\nobreak\hbox{-}}
\setlength{\emergencystretch}{2em}
\multlinegap0pt
\usepackage{amssymb}
\newtheorem{lemma}{Lemma}
\newcommand{\RR}{\mathbb{R}}
\def\bft{\mathbf{t}}
\title{A divided difference identity \\ for a class of multiple integrals}
\author{Michael S. Floater}
\date{Source: arXiv:2603.18867v1 (CC BY 4.0)}
\begin{document}
\maketitle

\noindent\emph{Source and licence.} A divided difference identity \\ for a class of multiple integrals. Version arXiv:2603.18867v1 (CC BY 4.0), distributed by arXiv under the Creative Commons Attribution 4.0 International licence, http://creativecommons.org/licenses/by/4.0/, as recorded in the arXiv licence metadata for this exact version and shown on the abstract page https://arxiv.org/abs/2603.18867v1. Full licence notice: \texttt{licenses/arxiv-2603.18867-CC-BY-4.0.txt}.\par\smallskip

\noindent\textit{Editorial scope.} Counted unit: the article's lemma lem:chain (source0.tex lines 489--499). The article's complete original proof of it is delivered with it (source0.tex lines 501--540, one \texttt{proof} environment of 40 lines). No mathematics is written, repaired or re-derived: the statement and the proof are the article's own lines, in the article's own notation, and no retained line is substituted or re-flowed.  No further source-local context is needed: every cross-reference of every retained line resolves inside the two delivered spans.  Nothing was omitted from any retained span, so the package carries no omission record.  No retained line contains a citation, so the wrapper carries no bibliography and no bibliography entry.  No retained line contains a figure environment, an image-inclusion command or drawing code, so no image is delivered and the package declares no asset dependency.  Editorial formatting only, in addition: an AMS \texttt{amsart} preamble that restates the article's own declarations and macros used by the retained lines, namely the environments \texttt{lemma} on separate counters, together with the standard packages those lines need.  The retained environments print in this document as Lemma 1 in order of appearance, whereas the article prints the same items with the numbering of its own full document.  The complete native source of the article as distributed in the arXiv e-print for this exact version is bundled unchanged as \texttt{sources/196731/source0.tex}.
\begin{lemma}\label{lem:chain}
Suppose that $\psi \in C^n(\Omega)$ and $f \in C^n(s(\Omega))$
for some domain $\Omega \subseteq \RR^n$, and define the product
$$ \phi(\bft) := \psi(\bft)f\big(s(\bft)\big), \qquad \bft \in \Omega. $$
Then,
\begin{equation}\label{eq:Echain}
\frac{\partial^n}{\partial t_1 \cdots \partial t_n} \phi(\bft)
= E_n\phi(\bft)
= \sum_{k=0}^n E_k \psi(\bft) f^{(n-k)}\big(s(\bft)\big), \quad \bft \in \Omega.
\end{equation}
\end{lemma}

\begin{proof}
We show more generally that
\begin{equation}\label{eq:Echainm}
\frac{\partial^m}{\partial t_1 \cdots \partial t_m} \phi(\bft)
= \sum_{k=0}^m
\sum_{1 \le i_1 < \cdots < i_k \le m}
\frac{\partial^k}{\partial t_{i_1} \cdots \partial t_{i_k}}
\psi(\bft) f^{(m-k)}\big(s(\bft)\big)
\end{equation}
for all $m=0,1,\ldots,n$ by induction on $m$.
This is true for $m=0$ by the definition of $\phi$.
Suppose it holds for some $m$, $0 \le m \le n-1$.
Differentiation with respect to $t_{m+1}$ gives
$$ 
\frac{\partial^{m+1}}{\partial t_1 \cdots \partial t_{m+1}} \phi(\bft)
   = a + b, $$
where
\begin{align*}
a &:= \sum_{k=0}^m 
\sum_{1 \le i_1 < \cdots < i_k \le m}
\frac{\partial^{k+1}}{\partial t_{i_1} \cdots \partial t_{i_k} \partial t_{m+1}}
\psi(\bft) f^{(m-k)}\big(s(\bft)\big), \cr
b &:= \sum_{k=0}^m 
\sum_{1 \le i_1 < \cdots < i_k \le m}
\frac{\partial^k}{\partial t_{i_1} \cdots \partial t_{i_k}}
\psi(\bft) f^{(m+1-k)}\big(s(\bft)\big).
\end{align*}
We can write $a$ as
\begin{align*}
a &= \sum_{k=1}^{m+1} 
\sum_{1 \le i_1 < \cdots < i_{k-1} \le m}
\frac{\partial^k}{\partial t_{i_1} \cdots \partial t_{i_{k-1}} \partial t_{m+1}}
\psi(\bft) f^{(m+1-k)}\big(s(\bft)\big) \cr
&= \sum_{k=1}^{m+1} 
\sum_{1 \le i_1 < \cdots < i_k = m+1}
\frac{\partial^k}{\partial t_{i_1} \cdots \partial t_{i_k}}
\psi(\bft) f^{(m+1-k)}\big(s(\bft)\big),
\end{align*}
and adding this to $b$ yields the case $m+1$ of (\ref{eq:Echainm}).
\end{proof}

\end{document}
